\section[An infinite-dimensional $\Re$-module and its universal property]{An infinite-dimensional $\boldsymbol{\Re}$-module and its universal property}\label{s:Verma}

In this section we introduce an infinite-dimensional $\Re$-module and its universal property in order to prove Theorem~\ref{thm:classification}. For convenience the following conventions are used throughout the rest of this paper. We let $a$, $b$, $c$, $\nu$ denote any scalars in~$\F$. We define the following families of parameters associated with $a$, $b$, $c$, $\nu$:
\begin{gather}
\theta_i =\big(a+\textstyle\frac{\nu}{2}-i\big)\big(a+\textstyle\frac{\nu}{2}-i+1\big)\qquad \hbox{for all $i\in \Z$}, \label{theta_i}\\
\theta^*_i=\big(b+\textstyle\frac{\nu}{2}-i\big)\big(b+\textstyle\frac{\nu}{2}-i+1\big)\qquad\hbox{for all $i\in \Z$}, \label{theta_is}\\
\phi_i=i(i-\nu-1)\big(a-b+c-\textstyle\frac{\nu}{2}+i\big)\big(a-b-c-\textstyle\frac{\nu}{2}+i-1\big)\qquad
\hbox{for all $i\in \Z$},\label{phi_i}\\
\varphi_i=i(i-\nu-1)\big(a+b+c+\textstyle\frac{\nu}{2}-i+2\big)\big(a+b-c+\textstyle\frac{\nu}{2}-i+1\big)
\qquad\hbox{for all $i\in \Z$},\label{varphi_i}\\
\zeta=(c-b)(c+b+1)\big(a-\textstyle\frac{\nu}{2}\big)\big(a+\textstyle\frac{\nu}{2}+1\big),\label{omega}\\
\zeta^*=(a-c)(a+c+1)\big(b-\textstyle\frac{\nu}{2}\big)\big(b+\textstyle\frac{\nu}{2}+1\big),\label{omages}\\
\eta=\textstyle\frac{\nu}{2} \big(\textstyle\frac{\nu}{2}+1\big)+a(a+1)+b(b+1)+c(c+1).\label{eta}
\end{gather}

\begin{Proposition}\label{prop:Verma}There exists an $\Re$-module $M_\nu (a,b,c)$ satisfying each of the following conditions:
\begin{enumerate}\itemsep=0pt
\item[{\rm (i)}] There exists an $\F$-basis $\{m_i\}_{i=0}^\infty$ for $M_\nu (a,b,c)$ with respect to which the matrices representing $A$ and $B$ are
\begin{gather*}
\begin{pmatrix}
\theta_0 & & & & &{\bf 0}
\\
1 &\theta_1
\\
&1 &\theta_2
 \\
 & &\cdot &\cdot
 \\
& & &\cdot &\cdot
 \\
 {\bf 0} & & & &\cdot &\cdot
\end{pmatrix},
\qquad
\begin{pmatrix}
\theta_0^* &\varphi_1 & & & &{\bf 0}
\\
 &\theta_1^* &\varphi_2
\\
 & &\theta_2^* &\cdot
 \\
 & & &\cdot &\cdot
 \\
& & & &\cdot &\cdot
 \\
 {\bf 0} & & & & &\cdot
\end{pmatrix},
\end{gather*}
respectively.

\item[{\rm (ii)}] The elements $\alpha$, $\beta$, $\delta$ act on $M_\nu (a,b,c)$ as scalar multiplication by $\zeta$, $\zeta^*$, $\eta$, respectively.
\end{enumerate}
\end{Proposition}
\begin{proof}Using Lemma \ref{lem:presentationR}, this result can be verified through routine computations.
\end{proof}

Throughout the rest of this paper we will let $\{m_i\}_{i=0}^\infty$ denote the $\F$-basis for $M_\nu(a,b,c)$ from Proposition~\ref{prop:Verma}(i). The following result is an immediate consequence of Proposition~\ref{prop:Verma}(i).

\begin{Lemma}\label{lem:mi}$m_{j+1}=\prod\limits_{h=i}^j (A-\theta_h)m_i$ for any $i,j\in \N$ with $ i\leq j$.
\end{Lemma}

Shortly we will describe the $\Re$-module $M_\nu(a,b,c)$ in an alternate way. To aid us in this endeavor, we first recall a Poincar\'{e}--Birkhoff--Witt basis for $\Re$.

\begin{Lemma}[{\cite[Theorem 5.1]{SH:2017-1}}]\label{lem:basisURA}
The elements
\begin{gather*}
A^i D^j B^k \alpha^r \delta^s \beta^t\qquad \hbox{for all $i,j,k,r,s,t\in \N$}\label{eq:basisURA}
\end{gather*}
form an $\F$-basis of~$\Re$.
\end{Lemma}

Let $I_\nu(a,b,c)$ denote the left ideal of $\Re$ generated by the elements
\begin{gather}
B-\theta_0^*,\label{U1}\\
(B-\theta_1^*)(A-\theta_0)-\varphi_1,\label{U2}\\
\alpha-\zeta,\qquad\beta-\zeta^*,\qquad\delta-\eta.
\label{U3}
\end{gather}
We now consider certain cosets of $\Re/I_\nu(a,b,c)$.

\begin{Lemma}\label{lem:cosets}For each $n\in \N$, each of the following holds:
\begin{enumerate}\itemsep=0pt
\item[$(i)$] $BA^n+I_\nu(a,b,c)$ is an $\F$-linear combination of $A^i+I_\nu(a,b,c)$ for all $0\leq i\leq n$,
\item[$(ii)$] $DA^n+I_\nu(a,b,c)$ is an $\F$-linear combination of $A^i+I_\nu(a,b,c)$ for all $0\leq i\leq n+1$,
\item[$(iii)$] $D^n+I_\nu(a,b,c)$ is an $\F$-linear combination of $A^i+I_\nu(a,b,c)$ for all $0\leq i\leq n$.
\end{enumerate}
\end{Lemma}
\begin{proof} (i) We proceed by induction on $n$. Since $I_\nu(a,b,c)$ contains the element (\ref{U1}), the statement holds for $n=0$. Since $I_\nu(a,b,c)$ contains both of the elements~(\ref{U1}) and~(\ref{U2}), the statement holds for $n=1$. Now suppose $n\geq 2$. Right multiplying each side of (\ref{AAB}) by $A^{n-2}$ yields that
\begin{gather*}
A^2 B A^{n-2}
-2 A B A^{n-1}
+B A^n
-2 A B A^{n-2}
-2 B A^{n-1}
=2 A^n-2 A^{n-1}\delta+2A^{n-2}\alpha.
\end{gather*}
Since $I_\nu(a,b,c)$ contains each of the elements listed in~(\ref{U3}), it follows that $B A^n$ is congruent to
\begin{gather}\label{ABAn-1}
2 A B A^{n-1}+2 A B A^{n-2}-A^2 B A^{n-2}+2 B A^{n-1}+2 A^n-2\eta A^{n-1}+2\zeta A^{n-2}
\end{gather}
modulo $I_\nu(a,b,c)$. By the inductive hypothesis, the element (\ref{ABAn-1}) is congruent to an $\F$-linear combination of $A^i$, for all $0\leq i\leq n$, modulo $I_\nu(a,b,c)$. Therefore (i) follows.

(ii) Observe that $DA^n=\frac{1}{2}\big(ABA^{n}-BA^{n+1}\big)$ by (\ref{r:D}). In light of this fact, the result now follows from Lemma~\ref{lem:cosets}(i).

(iii) We proceed by induction on $n$. The statement holds trivially for $n=0$. Now suppose that $n\geq 1$. By the inductive hypothesis, $D^n=DD^{n-1}$ is congruent to an $\F$-linear combination of
\begin{gather}\label{DA^i+I}
DA^i, \qquad 0\leq i\leq n-1,
\end{gather}
modulo $I_\nu(a,b,c)$. By Lemma \ref{lem:cosets}(ii) each of the elements listed in~(\ref{DA^i+I}) is an $\F$-linear combination of~$A^k$, for all $0\leq k\leq n$, modulo $I_\nu(a,b,c)$. Therefore the result follows.
\end{proof}

\begin{Lemma}\label{lem:R/I_basis}The $\F$-vector space $\Re/I_\nu(a,b,c)$ is spanned by
\begin{gather*}
A^i+I_\nu(a,b,c) \qquad \hbox{for all $i\in \N$}.
\end{gather*}
\end{Lemma}
\begin{proof}By Lemma \ref{lem:basisURA}, the $\F$-vector space $\Re/I_\nu(a,b,c)$ is spanned by
\begin{gather}\label{coset1}
A^i D^j B^k \alpha^r \delta^s \beta^t+I_\nu(a,b,c)\qquad \hbox{for all $i,j,k,r,s,t\in \N$}.
\end{gather}
Since $I_\nu(a,b,c)$ contains the elements listed in (\ref{U1}) and (\ref{U3}), each of the elements listed in~(\ref{coset1}) can be expressed as an $\F$-linear combination of
\begin{gather*}%\label{coset2}
A^i D^j+I_\nu(a,b,c) \qquad \hbox{for all $i,j\in \N$}.
\end{gather*}
The result now follows from these facts along with Lemma~\ref{lem:cosets}(iii).
\end{proof}

We are now ready to give our second description of $M_\nu(a,b,c)$.

\begin{Theorem}\label{thm:M=R/I}There exists a unique $\Re$-module homomorphism
\begin{gather*}
\Phi\colon \ \Re/I_\nu(a,b,c)\to M_\nu(a,b,c)
\end{gather*}
that sends $1+I_\nu(a,b,c)$ to $m_0$. Moreover, $\Phi$ is an isomorphism.
\end{Theorem}
\begin{proof}Consider the $\Re$-module homomorphism $\Psi\colon \Re\to M_\nu(a,b,c)$ that sends $1$ to $m_0$. By Proposition \ref{prop:Verma}(i), the elements (\ref{U1}) and (\ref{U2}) are in the kernel of $\Psi$. By Proposition \ref{prop:Verma}(ii), the elements listed in (\ref{U3}) are also in the kernel of $\Psi$. Hence $I_\nu(a,b,c)$ is contained in the kernel of~$\Psi$. It follows that $\Psi$ induces an $\Re$-module homomorphism $\Phi\colon \Re/I_\nu(a,b,c) \to M_\nu(a,b,c)$ that maps $1+I_\nu(a,b,c)$ to~$m_0$. Observe that $\Phi$ is the unique $\Re$-module homomorphism with the desired property since $\Re/I_\nu(a,b,c)$ is generated by $1+I_\nu(a,b,c)$ as an $\Re$-module.

By Lemma \ref{lem:mi} the homomorphism $\Phi$ sends
\begin{gather}\label{coset3}
\prod_{h=1}^{i-1}(A-\theta_h)+I_\nu(a,b,c)
\end{gather}
to $m_i$ for all $i\in \N$. Since $\{m_i\}_{i=0}^\infty$ are linearly independent, the cosets~(\ref{coset3}) are linearly independent. Combining this with Lemma~\ref{lem:R/I_basis}, we see that the cosets~(\ref{coset3}) are an $\F$-basis for $\Re/I_\nu(a,b,c)$. Therefore $\Phi$ is an isomorphism.
\end{proof}

As a consequence of Theorem~\ref{thm:M=R/I}, the $\Re$-module $M_\nu(a,b,c)$ satisfies the following universal property.

\begin{Proposition}\label{prop:universal}If $V$ is an $\Re$-module which has a vector $v\in V$ satisfying
\begin{gather*}
Bv=\theta_0^* v,\\
(B-\theta_1^*)(A-\theta_0)v=\varphi_1 v,\\
\alpha v=\zeta v,\qquad\beta v=\zeta^* v,\qquad\delta v=\eta v,
\end{gather*}
then there exists a unique $\Re$-module homomorphism $M_\nu(a,b,c)\to V$ that sends~$m_0$ to~$v$.
\end{Proposition}

For the rest of the present paper, we will consider the case $\nu=d$. Define $N_d(a,b,c)$ to be the $A$-cyclic $\F$-subspace of $M_d(a,b,c)$ generated by the element $m_{d+1}$.

\begin{Lemma}\label{lem:N}$N_d(a,b,c)$ is an $\Re$-submodule of $M_d(a,b,c)$ with the $\F$-basis $\{m_i\}_{i=d+1}^\infty$.
\end{Lemma}
\begin{proof}Recall from Lemma \ref{lem:mi} that
\begin{gather*}
(A-\theta_i) m_i=m_{i+1}\qquad \hbox{for all $i\geq d+1$}.
\end{gather*}
It follows from this fact that $\{m_i\}_{i=d+1}^\infty$ is an $\F$-basis for $N_d(a,b,c)$.

We now show that $N_d(a,b,c)$ is an $\Re$-submodule of $M_d(a,b,c)$. By Proposition~\ref{prop:Verma}(i),
\begin{gather*}
(B-\theta_i^*)m_i=\varphi_i m_{i-1} \qquad\hbox{for all $i\geq d+1$}.
\end{gather*}
By (\ref{varphi_i}), the scalar $\varphi_{d+1}=0$ when $\nu=d$. Hence $N_d(a,b,c)$ is $B$-invariant. By Proposition~\ref{prop:Verma}(ii), the element $\delta$ acts on $N_d(a,b,c)$ as scalar multiplication by $\eta$.
It now follows from Lemma \ref{lem:delta}(ii) that
 $N_d(a,b,c)$ is an $\Re$-submodule of $M_d(a,b,c)$.
\end{proof}

Recall the $\Re$-module $R_d(a,b,c)$ from Proposition~\ref{prop:Rd}. In the sequel we display how the $\Re$-module $R_d(a,b,c)$ is connected to $M_\nu(a,b,c)$. For convenience we let $\{v_i\}_{i=0}^d$ denote the $\F$-basis for $R_d(a,b,c)$ from Proposition~\ref{prop:Rd}(i) in the rest of this paper.

\begin{Lemma}\label{lem:M/N}There exists a unique $\Re$-module isomorphism
\begin{gather*}
M_d(a,b,c)/N_d(a,b,c)\to R_d(a,b,c)
\end{gather*}
that sends $m_i+N_d(a,b,c)$ to $v_i$ for all $0\leq i\leq d$.
\end{Lemma}
\begin{proof}By Lemma \ref{lem:N}, $M_d(a,b,c)/N_d(a,b,c)$ is a $(d+1)$-dimensional $\Re$-module with the $\F$-basis
\begin{gather}\label{mi+N}
\{m_i+N_d(a,b,c)\}_{i=0}^d.
\end{gather}
Observe that the matrices representing $A$ and $B$ with respect to the $\F$-basis $\{v_i\}_{i=0}^d$ for $R_d(a,b,c)$ are identical with the matrices representing $A$ and $B$ with respect to the $\F$-basis (\ref{mi+N}) for $M_d(a,b,c)/N_d(a,b,c)$ by Propositions~\ref{prop:Rd}(i) and~\ref{prop:Verma}(i). By Propositions~\ref{prop:Rd}(ii) and~\ref{prop:Verma}(ii), the actions of $\delta$ on $R_d(a,b,c)$ and $M_d(a,b,c)/N_d(a,b,c)$ are scalar multiplication by the same scalar~$\eta$. In light of these comments, the result now follows from Lemma~\ref{lem:delta}(ii).
\end{proof}

\begin{Proposition}\label{prop:R}Suppose that $V$ is an $\Re$-module which has a vector $v\in V$ satisfying
\begin{gather}\label{md+1_ker}
\prod_{i=0}^d (A-\theta_i) v=0.
\end{gather}
If there is an $\Re$-module homomorphism $M_d(a,b,c)\to V$ that sends $m_0$ to~$v$, then there exists an $\Re$-module homomorphism $R_d(a,b,c)\to V$ that sends~$v_0$ to~$v$.
\end{Proposition}
\begin{proof}Let $\varrho$ denote the $\Re$-module homomorphism $M_d(a,b,c)\to V$ that sends~$m_0$ to~$v$. By Lemma~\ref{lem:mi}, we have
\begin{gather*}
m_{d+1}=\prod_{i=0}^d (A-\theta_i) m_0.
\end{gather*}
Combining this with~(\ref{md+1_ker}), we see that $m_{d+1}$ is in the kernel of $\varrho$. Therefore $N_d(a,b,c)$ is contained in the kernel of~$\varrho$. By Lemma~\ref{lem:N}, there exists an $\Re$-module homomorphism $M_d(a,b,c)/N_d(a,b,c)\to V$ that sends $m_0+N_d(a,b,c)$ to~$v$. The result follows from this fact along with Lemma~\ref{lem:M/N}.
\end{proof}
